\mysection{缩略词表}

% Tightened locally so the whole list still fits on one page. longtable (rather
% than tabular) is deliberate: if the list outgrows the page later it breaks
% across pages instead of silently running off the bottom.
\begingroup
\small
\renewcommand{\arraystretch}{1.0}
\setlength{\LTpre}{0pt}
\setlength{\LTpost}{0pt}
\begin{longtable}{@{}p{1.8cm} p{11.8cm}@{}}
\multicolumn{2}{@{}l}{\textbf{实验与物理}} \\[2pt]
CFD   & 计算流体力学（Computational Fluid Dynamics） \\
DBI   & 漫射背光成像（Diffuse Background Illumination） \\
ECN   & 发动机燃烧网络（Engine Combustion Network） \\
EOI   & 喷射结束（End of Injection） \\
FOV   & 视场（Field of View） \\
OSCC  & 光学喷雾燃烧舱（Optical Spray Combustion Chamber） \\
SOI   & 喷射起始（Start of Injection） \\[2pt]
\multicolumn{2}{@{}l}{\textbf{图像处理与计算}} \\[2pt]
BW    & 二值化图像（Binary Image） \\
CDF   & 累积分布函数（Cumulative Distribution Function） \\
CPU   & 中央处理器（Central Processing Unit） \\
CUDA  & 统一计算设备架构（Compute Unified Device Architecture） \\
FFT   & 快速傅里叶变换（Fast Fourier Transform） \\
FPS   & 帧率，帧/秒（Frames Per Second） \\
GPU   & 图形处理器（Graphics Processing Unit） \\
QC    & 质量控制（Quality Control） \\
SIMT  & 单指令多线程（Single Instruction Multiple Thread） \\[2pt]
\multicolumn{2}{@{}l}{\textbf{机器学习}} \\[2pt]
ARD   & 自动相关性确定（Automatic Relevance Determination） \\
FiLM  & 逐特征线性调制（Feature-wise Linear Modulation） \\
GP    & 高斯过程（Gaussian Process） \\
KD    & 知识蒸馏（Knowledge Distillation） \\
KL    & 库尔贝克--莱布勒散度（Kullback--Leibler Divergence） \\
LONO  & 留一喷嘴评估（Leave-One-Nozzle-Out） \\
MLP   & 多层感知机（Multilayer Perceptron） \\
NLL   & 负对数似然（Negative Log-Likelihood） \\
OOD   & 分布外（Out-of-Distribution） \\
SGQR  & Sigmoid 门控四次方根（Sigmoid-Gated Quarter-Root） \\
SVGP  & 稀疏变分高斯过程（Sparse Variational Gaussian Process） \\[2pt]
\multicolumn{2}{@{}l}{\textbf{统计与评估指标}} \\[2pt]
CRPS  & 连续排名概率分数（Continuous Ranked Probability Score） \\
ECE   & 期望校准误差（Expected Calibration Error） \\
MAD   & 中位数绝对偏差（Median Absolute Deviation） \\
MAE   & 平均绝对误差（Mean Absolute Error） \\
MSE   & 均方误差（Mean Squared Error） \\
PIT   & 概率积分变换（Probability Integral Transform） \\
RMSE  & 均方根误差（Root Mean Square Error） \\
SCR   & 空间删失率（Spatial Censoring Rate） \\
\end{longtable}
\endgroup
